\section{Literature Review}\label{sec:literature}

This section situates our work within five interconnected streams of research: (1) cognitive load and attention management, (2) job demands, resources, and burnout, (3) work-life balance and flexibility, (4) technology and AI in the workplace, and (5) optimization approaches to resource allocation.

\subsection{Cognitive Load and Attention Management}

Cognitive load theory provides foundational insights for sustainable life management. Research demonstrates that cognitive training yields lasting benefits, with reasoning and speed-of-processing training showing effects up to 10 years in older adults, supporting strategies for sustaining cognitive resources across the lifespan \citep{NIA2014}. Contemporary applications of cognitive load theory in computer-based environments reveal the importance of worked examples, split-attention effects, and the self-management of cognitive load \citep{ayres_something_2020}. NSF-funded research on peak productivity develops tools using smartphones and smartwatches to track biological clocks and align tasks with peak alertness periods, addressing disruptions from mismatched timing that cause health issues and errors \citep{Lefkowitz2018}.

The distinction between time management and attention management is critical. Traditional time management strategies fail because they ignore cognitive constraints, forcing individuals to confront limitations and make hard choices among competing priorities rather than attempting to optimize everything \citep{Hart2024}. Productivity hinges more on managing attention than rigid time scheduling, as cognitive limitations make it difficult to sustain focus across competing demands \citep{Grant2019}. This recognition has spurred interest in decision support systems that account for cognitive constraints. Research shows that just-in-time access to codified knowledge reduces cognitive load and enables learning, with workers accessing procedural knowledge from document repositories building firm-specific human capital \citep{subramani_capability_2021}.

Human-AI collaboration offers promise for reducing cognitive burden but requires careful design. Studies find that granting users decision control improves trust and compliance with AI recommendations, while explanations may increase task complexity depending on the user's cognitive ability \citep{westphal_decision_2023}. The challenge is particularly acute for knowledge workers facing information overload. Mitigation strategies include streamlining information, investing in training and mental health resources, promoting deep work with breaks, and ensuring equitable workload distribution \citep{Robinson2024}.

\subsection{Job Demands, Resources, and Burnout}

The Job Demands-Resources (JD-R) model provides a powerful framework for understanding sustainable life management. Research confirms that job resources moderate the relationship between job demands and occupational strain, with autonomy reducing strain under high workload conditions \citep{ghanayem_occupational_2020}. Extensions of JD-R theory beyond the work domain reveal that senior executives utilize both internal and external resources, and that job crafting creates adaptable resources to meet evolving demands \citep{duan_extending_2024}. Individual health status itself functions as a resource: chronic illness symptom severity predicts burnout and reduced work engagement above and beyond the effects of job demands and resources \citep{cook_individual_2024}.

Conservation of Resources (COR) theory complements JD-R by explaining resource dynamics. Research distinguishes emotional resource accumulation from depletion, developing measurement instruments that capture how resource-generating activities (work engagement, co-worker support) and resource-depleting activities (job demands, emotional labor) differentially predict family functioning and life satisfaction \citep{ilies_emotional_2020}. Supervisor resilience can cross over to employees, with time-lagged studies demonstrating that supervisor resilience promotes employee well-being through lower burnout and psychological distress \citep{brady_supervisor_2025}. Managers can help employees navigate resource-intensive tasks through equipping, conserving, and restoring the self across task phases \citep{molinsky_how_2023}.

Burnout represents a critical outcome when resource allocation fails. Healthcare providers face particular challenges: physicians' workflow decisions regarding when to perform electronic health record tasks significantly influence total workload and after-hours work, with implications for sustainable practice \citep{celik_frontiers_2024}. STEM women recovering from burnout employ distinct sensemaking strategies (combative, regenerative, and promissory) to renegotiate work trajectories, revealing the relational negotiation of protecting, investing, and fostering resources during disruption \citep{lee_beyond_2024}. Workplace interventions like the Live Healthy, Work Healthy program improve perceived organizational support, reduce burnout, and increase work engagement through structured support for chronic condition management \citep{haynes_evaluating_2021}.

\subsection{Work-Life Balance and Flexibility}

Work-family conflict research has evolved from composite-level to dimension-level analysis. Meta-analytic evidence demonstrates that time-based, strain-based, and behavior-based dimensions of conflict have distinct theoretical underpinnings and empirical correlates, with resource allocation theory particularly relevant to time-based conflict \citep{hetrick_theoretical_2023}. Crowdsourced firm ratings analysis reveals that work-life balance ratings influence organizational outcomes differently across industries, suggesting that decision makers should tailor motivation strategies to workforce needs \citep{liu_crowdsourced_2024}.

Flexible working arrangements (FWAs) offer promising solutions to work-life challenges. Comprehensive meta-analysis identifies antecedents of FWA availability and use (managerial status, self-efficacy, task interdependency) and confirms beneficial outcomes including job satisfaction, organizational commitment, and family satisfaction \citep{harrop_meta_2025}. Importantly, having both flextime and flexplace proves more beneficial than either alone, and availability may matter more than actual use. Experimental evidence from a randomized trial demonstrates that flexible place and time of work increases productivity while improving work-life balance, with men also increasing household and care time \citep{angelici_smart_2024}.

Alternative work schedules provide additional flexibility mechanisms. Longitudinal research on compressed workweek implementation shows that work-life balance increases while fatigue and time pressure decrease, with effects contingent on employee expectations regarding the schedule change \citep{mühl_you_2024}. The future of work increasingly involves hybrid arrangements despite return-to-office pressures, with research showing that work-from-anywhere policies can enhance performance through flexibility \citep{McGregor2023}.

\subsection{Technology and AI in the Workplace}

Technology plays a dual role in sustainable life management, serving as both enabler and stressor. Technostressors from information technology use have become a major stress source, negatively affecting work-life balance through emotional exhaustion, though job self-efficacy buffers these effects \citep{ma_impact_2021}. Configurational analysis reveals that interdependencies among technostressors (complementarity, contingency, and substitution) form patterns leading to burnout and reduced job performance, explaining why interventions addressing independent technostressors may fail \citep{pflügner_deconstructing_2024}. Information and communication technology can reduce nurses' burden by facilitating environmental management and non-contact communication while providing emotional support for patients, though implementation requires attention to digital health literacy and safety \citep{yoo_critical_2022}.

Artificial intelligence introduces new dynamics into resource allocation. Reliance on algorithmic management in healthcare risks depleting human cognitive resources by favoring data summaries over nuanced listening, fostering cognitive miserliness and deskilling \citep{Reinhart2025}. In the gig economy, algorithmic evaluation and discipline increase burnout by negatively impacting psychological needs, while algorithmic direction can alleviate burnout by enhancing autonomy \citep{peng_every_2025}. Research on augmentation-based AI reveals countervailing mechanisms: frequent usage leads to knowledge gain and improved performance, but also causes information overload impairing performance and recovery \citep{shao_using_2024}. Generative AI can serve as a helpful assistant for both strategic and operational tasks when implemented with appropriate prompts and verification processes \citep{aguinis_how_2024}. Predictions of near-total AI dominance in key domains by 2027 raise fundamental questions about how humans will allocate limited time and cognitive effort as work systems evolve \citep{Douthat2025}.

Human-robot collaboration in manufacturing demonstrates the potential for dynamic task allocation. Research on recyclofacturing emphasizes dynamic task distribution between humans and robots in assembly and welding, with attention to human factors for safe collaboration \citep{SCU2025}. The key insight is that technology should adapt task distribution to evolving conditions and recovery needs in human-machine teams.

\subsection{Optimization Approaches to Resource Allocation}

Operations research provides powerful tools for sustainable life management. Constraint programming models for personal process management decrease cognitive load from decision-making and generate positive user experiences, enabling better planning in less time and fast replanning in dynamically changing situations \citep{oruç_constraint_2022}. Behavioral research demonstrates that simpler yet efficient contracts can mitigate the trade-off between efficiency and complexity, explaining the prevalence of minimum order quantity terms in supply contracts \citep{tuncel_why_2022}. Understanding incentive structures is critical: research shows that incentivizing solely estimation accuracy can result in hidden inefficiency from strategic schedule inflation, while combined accuracy and speed incentives yield compressed but reliable schedules \citep{lorko_hidden_2023}.

Healthcare resource allocation presents particularly challenging optimization problems. First Route Second Assign decomposition approaches enforce continuity of care in home health care while improving routing efficiency \citep{lahrichi_first_2022}. Two-stage stochastic game models guide capacity planning investments to address surgical waiting lists while accounting for demand uncertainty \citep{acuna_two_2024}. Multistage staffing strategies utilizing capacities of various flexibility help manage service systems facing worker shortages, preventing burnout in understaffed systems while improving service quality \citep{yuan_managing_2025}.

Workforce sizing and scheduling optimization can balance multiple objectives. Research on last-mile delivery demonstrates that introducing limited shift flexibility achieves similar effects as completely flexible systems while ensuring better work-life balance for workers, showing that improved working conditions need not sacrifice profitability \citep{mandal_tactical_2025}. This insight that stability and efficiency can coexist underlies our approach to sustainable life management.

\subsection{Research Gap and Contribution}

Despite the extensive literature on individual aspects of cognitive load, burnout, work-life balance, technology, and optimization, no existing framework integrates these streams into a unified optimization model for personal resource allocation. Prior work either focuses on organizational interventions without mathematical formalization or develops optimization models without accounting for cognitive constraints and recovery dynamics. Our contribution addresses this gap by developing a tractable mixed-integer stochastic programming formulation that captures state-dependent prioritization, uncertain recovery requirements, and multi-dimensional value creation, providing a foundation for decision support systems that promote sustainable life management.
